\documentclass[uplatex,dvipdfmx]{jsarticle}
\usepackage{amsmath,amssymb,amsthm}
\usepackage{enumitem}
\usepackage{seqsplit}
\usepackage[
  dvipdfmx,
  unicode,
  bookmarks=true,
  bookmarksnumbered=true,
  bookmarksopen=true,
  bookmarksdepth=2
]{hyperref}
\usepackage{bookmark}
\hypersetup{
  hidelinks,
  pdftitle={A Witt Local--Global Principle over a Rational Function Field},
  pdfauthor={Mocho Go},
  pdfsubject={Complete mathematical proof corresponding to the Lean formalization}
}

\newcommand{\RF}{R(X)}
\newcommand{\RatFunc}{\operatorname{RatFunc}}
\newcommand{\ord}{\operatorname{ord}}
\newcommand{\sgn}{\operatorname{sgn}}
\newcommand{\lean}[1]{\texttt{\expandafter\seqsplit\expandafter{\detokenize{#1}}}}

\newtheorem{theorem}{Theorem}
\newtheorem{proposition}[theorem]{Proposition}
\newtheorem{lemma}[theorem]{Lemma}
\newtheorem{corollary}[theorem]{Corollary}
\theoremstyle{definition}
\newtheorem{definition}[theorem]{Definition}
\theoremstyle{remark}
\newtheorem{remark}[theorem]{Remark}
\renewcommand{\abstractname}{Abstract}

\title{A Witt Local--Global Principle over a Rational Function Field}
\author{Mocho Go}
\date{}

\begin{document}
\pdfbookmark[0]{A Witt Local--Global Principle over a Rational Function Field}{document-title}
\maketitle

\begin{abstract}
Let $R$ be a real closed field and $F=R(X)$.  We give the proof formalized
in the Lean project \texttt{lean-ratfunc-witt-local-global}: every finite-dimensional
quadratic form over $F$, of dimension at least three, is isotropic if and
only if it is isotropic after extension to every real closed ordered extension
of $F$.  The proof constructs an ordered real closure for the ternary base
case and then uses a sign-separating binary--tail compression to induct on the
dimension.  We also record the correspondence between the mathematical lemmas
and the public Lean declarations.
\end{abstract}

\section{Statement and conventions}

Throughout, $R$ is a real closed field and $F=R(X)$.  If $I$ is finite
and $a=(a_i)_{i\in I}$ is a family in $F$, write
\[
  \langle a\rangle(x)=\sum_{i\in I}a_i x_i^2.
\]
The diagonal form is \emph{isotropic} if this sum vanishes for some nonzero
vector $x\in F^I$.  An ordering $P$ of $F$ makes a regular diagonal form
definite if all coefficients are strictly positive in $P$, or all are
strictly negative in $P$.

\begin{theorem}[Diagonal local--global theorem]\label{thm:diagonal}
Let $I$ be finite with $|I|\ge3$, and let $a:I\to F$.  The following are
equivalent.
\begin{enumerate}[label=(\roman*)]
  \item $\langle a\rangle$ is isotropic over $F$.
  \item No ordering of $F$ makes all coefficients of $a$ have the same
        strict sign.
  \item For every real closed ordered field $K$ equipped with an
        $F$-algebra structure, the scalar extension
        $\langle a\rangle_K$ is isotropic.
\end{enumerate}
In (iii), the Lean theorem quantifies over target fields in a fixed universe;
for another universe it takes the existence of a corresponding ordered real
closure as an explicit input.
\end{theorem}

\begin{theorem}[Quadratic-form local--global theorem]\label{thm:qf}
Let $V$ be a finite-dimensional $F$-vector space with
$\dim_F V\ge3$, and let $q:V\to F$ be a quadratic form.  Then $q$ is not
anisotropic if and only if every scalar extension of $q$ to a real closed
ordered extension of $F$ is not anisotropic.
As in Theorem~\ref{thm:diagonal}, the formal statement quantifies over
same-universe targets; a universe-parameterized version assumes a suitable
ordered real closure.
\end{theorem}

Zero coefficients cause no difficulty: a coordinate vector is immediately
isotropic.  Thus the proof may assume all coefficients are nonzero until the
final total statement.

\paragraph{Reading guide.}
The proof uses fields, unique
factorization in $R[X]$, the Chinese remainder theorem, finite-dimensional
linear algebra, and Zorn's lemma.  The less standard ingredients are developed
where they are first needed.  Section~2 constructs the ordered real closure
used to turn orderings into genuine test fields.  Section~3 proves the ternary
case by translating a diagonal form into a quaternion conic and carrying out
the same polynomial descent as the Lean development.  Sections~4--5 construct
one rational function with prescribed signs and use it to reduce the dimension
by one.  Section~6 passes from diagonal forms to arbitrary quadratic forms.
The final table gives the Lean declaration implementing every major step, so a
reader may move in either direction between this proof and the source.

\section{Orderings and real closed extensions}

We first isolate two standard observations.

\begin{lemma}[Definiteness obstruction]\label{lem:definite}
If a diagonal form is isotropic over an ordered field, its coefficients
are neither all strictly positive nor all strictly negative.
\end{lemma}
\begin{proof}
If all coefficients are positive, every nonzero summand $a_i x_i^2$ is
nonnegative and at least one is positive.  Hence their sum is positive.
Negating the form treats the all-negative case.
\end{proof}

\begin{lemma}[Diagonal forms over a real closed field]\label{lem:rcdiag}
Let $K$ be real closed and let all coefficients $b_i\in K$ be nonzero.
Then $\langle b\rangle$ is anisotropic if and only if all $b_i$ have the
same strict sign.
\end{lemma}
\begin{proof}
The reverse implication is Lemma~\ref{lem:definite}.  Otherwise choose
the inequalities $b_i>0>b_j$.  Since $-b_j/b_i>0$, it is a square in $K$; if
$u^2=-b_j/b_i$, the vector with coordinate $u$ at $i$, coordinate $1$ at $j$,
and zero elsewhere is isotropic.
\end{proof}

The formalization does not assume an ordered real closure as an extra axiom.
It constructs one in the same universe.  We give the construction in some
detail, including the polynomial induction used for odd extensions, because
this is the input that makes the ternary statement unconditional.

\subsection{Two ordering-extension lemmas}

A \emph{preordering} of a field $L$ is a subset containing $0$ and every
square and closed under addition and multiplication.  It is \emph{proper} if
$-1$ does not belong to it.  An ordering is the same as a proper preordering
$P$ for which $P\cup(-P)=L$.

\begin{lemma}[Artin--Schreier extension lemma]\label{lem:preorder-extension}
Every proper preordering of a field is contained in an ordering.
\end{lemma}
\begin{proof}
Order the proper preorderings containing the given one by inclusion.  The
union of a chain is again a proper preordering, so Zorn's lemma gives a maximal
one, say $P$.  We show that it is total.  If neither $x$ nor $-x$ belongs to
$P$, maximality says that the preorderings generated by $P\cup\{x\}$ and by
$P\cup\{-x\}$ both contain $-1$.  Since powers of $x$ may be separated into
even and odd powers, these two statements have the form
\[
  -1=p+xq,\qquad -1=r-xs
  \quad(p,q,r,s\in P).
\]
Thus $xq=-(1+p)$ and $xs=1+r$.  Multiplication gives
\[
  x^2qs=-(1+p)(1+r)\in P.
\]
The left side belongs to $P$, so the displayed negative element does too.
Adding $p+r+pr\in P$ gives $-1\in P$, a contradiction.
Thus $x\in P$ or $-x\in P$ for every $x$, and $P$ is an ordering.
\end{proof}

\begin{lemma}[Positive projection criterion]\label{lem:projection}
Let $L/E$ be a field extension of an ordered field.  Suppose that there is an
$E$-linear map $\lambda:L\to E$ such that
\[
  z\ne0\quad\Longrightarrow\quad \lambda(z^2)>0.
\]
Then the ordering of $E$ extends to an ordering of $L$.
\end{lemma}
\begin{proof}
Let $P$ be the preordering of $L$ generated by the nonnegative elements of
$E$ and by the squares of $L$.  Every element of $P$ is a finite sum
$\sum_i a_i z_i^2$ with $a_i\ge0$.  If this sum were $-1$, applying
$\lambda$ would make its right-hand side nonnegative.  On the other hand,
the hypothesis applied to $1\ne0$ gives $\lambda(1)>0$, and hence
\[
  \lambda(-1)=-\lambda(1)<0.
\]
This is impossible.  Hence $P$ is proper, and
Lemma~\ref{lem:preorder-extension} applies.  The resulting cone contains the
old nonnegative elements.  If $a>0$ in $E$ but $-a$ also belonged to the new
cone, then multiplication by the old positive element $a^{-1}$ would put
$-1$ in that cone.  Thus old positive elements remain strictly positive, so
the old ordering is indeed extended.
\end{proof}

\subsection{Odd power-basis extensions}

The odd-degree part uses the following elementary polynomial lemma.  It is
included because the parity argument is easy to lose if it is compressed into
the phrase ``an odd extension is orderable.''

\begin{lemma}[Odd power-basis lemma]\label{lem:odd-power-basis}
Let $E$ be ordered and let $L=E(\theta)$ have a power basis of odd dimension
$n$.  Then the preordering generated by the nonnegative elements of $E$ and
the squares of $L$ is proper.
\end{lemma}
\begin{proof}
Let $f$ be the monic irreducible polynomial of $\theta$, so $\deg f=n$.
For $m\ge1$, let $\mathcal C_m$ be the set of polynomials that are finite sums
\[
  g=\sum_\nu a_\nu h_\nu^2,\qquad
  a_\nu\ge0,\quad \deg h_\nu<m.
\]
If $g\ne0$ lies in $\mathcal C_m$, the terms of largest degree cannot cancel:
their leading coefficients are sums of nonnegative squares.  Consequently
\[
  \operatorname{lc}(g)>0,\qquad
  \deg g\ \hbox{is even},\qquad
  \deg g<2m.                                      \tag{1}
\]

We prove by strong induction on the odd number $m$ that, for every monic
irreducible $r$ of degree $m$ and every $g\in\mathcal C_m$,
\[
  r\nmid g+1.                                      \tag{2}
\]
Suppose instead that $g+1=rk$.  The polynomial $g+1$ is nonzero because a
member of $\mathcal C_m$ cannot equal $-1$.  Since $r$ has positive
degree and divides $g+1$, the latter cannot be constant.  Thus $g$ is
nonconstant and $\deg(g+1)=\deg g$.  By (1),
\[
  \deg g=m+\deg k<2m.
\]
Since $\deg g$ is even and $m$ is odd, $\deg k$ is odd; moreover
$\deg k<m$.  Factorization in $E[T]$ therefore supplies a monic irreducible
factor $r'$ of $k$ having odd degree.  In
$E'=E[T]/(r')$, evaluation sends $k$ to zero, hence sends $g$ to $-1$.
Evaluation also sends $\mathcal C_m$ into the additive span of nonnegative
$E$-multiples of squares in $E'$.

Every element of $E'$ has a representative of degree $<\deg r'$.  Replacing
each square by such a representative lifts the equality $g(\bar T)=-1$ to a
polynomial $g'\in\mathcal C_{\deg r'}$ with
$g'(\bar T)=-1$.  Thus $r'\mid g'+1$, contrary to the induction hypothesis
because $\deg r'<m$.  This proves (2).

If the preordering in $L$ contained $-1$, writing its finitely many square
terms in the power basis would produce $g\in\mathcal C_n$ with
$g(\theta)=-1$.  The minimal polynomial $f$ would divide $g+1$, contradicting
(2).
\end{proof}

\begin{proposition}[Ordered real closure]\label{prop:orc}
Every ordered field $E$ admits an algebraic real closed ordered extension
whose order extends that of $E$.
\end{proposition}
\begin{proof}
Fix an algebraic closure $\Omega/E$.  Consider intermediate fields
$M\subseteq\Omega$ equipped with an ordering extending that of $E$, ordered
by inclusion with compatible positive cones.  Unions of chains again carry a
compatible ordering, so Zorn's lemma gives a maximal $M$.

We verify the two conditions in the ordered-field characterization of real
closed fields.

First let $0<a\in M$.  Choose $\alpha\in\Omega$ with $\alpha^2=a$.  If
$a$ is not already a square in $M$, then $T^2-a$ is irreducible and
$L=M(\alpha)$ is a proper quadratic extension.  In coordinates
$z=x+y\alpha$, the $M$-linear projection to the real coordinate satisfies
\[
  \pi(z^2)=x^2+a y^2>0\qquad(z\ne0).
\]
Lemma~\ref{lem:projection} orders $L$ over $M$, contradicting maximality.
Hence every nonnegative element of $M$ is a square.

Second, let $p\in M[T]$ have odd degree.  It has a monic irreducible factor
$r$ of odd degree: in any factorization of $p$, the sum of the degrees is
odd, so one irreducible factor has odd degree.  If $\alpha\in\Omega$ is a
root of $r$, then $M(\alpha)/M$ has an odd-dimensional power basis.
Lemma~\ref{lem:odd-power-basis} and
Lemma~\ref{lem:preorder-extension} order $M(\alpha)$ over $M$.
Maximality forces $\alpha\in M$, so every odd-degree polynomial has a root
in $M$.

An ordered field in which every nonnegative element is a square and every
odd-degree polynomial has a root is real closed.  Therefore $M$ is the
required extension.
\end{proof}

\section{The ternary base theorem}

The induction begins with the following theorem.

\begin{theorem}[Ternary local--global theorem]\label{thm:ternary}
Let $a,b,c\in F^\times$.  Then
$\langle a,b,c\rangle$ is isotropic over $F$ if and only if it is
isotropic over every real closed ordered extension of $F$.
\end{theorem}

We now give the algebraic core actually used in the Lean proof.  It is a
polynomial Legendre argument organized as a quaternion-symbol descent.

\subsection{Quaternion symbols and polynomial normalization}

For $u,v\in F^\times$, write $(u,v)$ for the quaternion symbol.  We shall use
only its elementary conic interpretation:
\[
 (u,v)\ \hbox{splits}
 \quad\Longleftrightarrow\quad
 \exists (x,y,z)\ne(0,0,0),\quad x^2-u y^2-v z^2=0.       \tag{3}
\]
Thus
\[
 \langle a,b,c\rangle\ \hbox{is isotropic}
 \quad\Longleftrightarrow\quad
 \left(-\frac ba,-\frac ca\right)\ \hbox{splits}.         \tag{4}
\]
Multiplying either slot by a nonzero square does not change splitting.
We also use the norm move
\[
 (u,v)\simeq (u,v(x^2-u y^2)),                            \tag{5}
\]
whenever the displayed multiplier is nonzero.  Formula (5) follows directly
from
\[
 (x^2-u y^2)(X^2-uY^2)
   =(xX+u yY)^2-u(yX+xY)^2.
\]
Thus multiplication by the norm $x^2-u y^2$ is an invertible change in the
binary norm equation.  If the third coordinate in (3) is nonzero, divide by
its square and apply this identity; if it is zero, $u$ is already a square
and both symbols split.  Applying the same calculation to the inverse norm
gives the converse.
Together with symmetry and cyclic permutation of a ternary form, these are
all quaternion identities used below.

Clearing denominators, multiplying the whole ternary form by one common
nonzero scalar, and removing an individual square from every coefficient
gives nonzero squarefree polynomials $A,B,C$.  They may be chosen pairwise
coprime.  Indeed, at each irreducible polynomial $\pi$, only the parities of
the three valuations matter.  Individual square removal makes the parities
$0$ or $1$, and multiplying all three coefficients by $\pi$ flips all three
bits.  We choose the representative in which at most one bit is $1$.
These operations preserve isotropy and preserve the property that all three
coefficients have one strict sign at an ordering.

For such a triple put
\[
  \mathcal Q(A,B,C)=(-AB,-AC).
\]
It differs from the normalized symbol in (4) only by square factors.

\subsection{Finite congruences}

We realize signs at a finite point and at infinity by Laurent series.  The
ordered field $R((\varepsilon))$ is ordered by declaring a nonzero series
positive when its first nonzero coefficient is positive.  The embeddings
\[
  X\longmapsto r+\varepsilon,\qquad
  X\longmapsto r-\varepsilon
\]
give the orderings immediately to the right and left of $r$; the embeddings
$X\mapsto\varepsilon^{-1}$ and
$X\mapsto-\varepsilon^{-1}$ give the two orderings at infinity.  A polynomial
with a simple root at $r$ has leading Laurent term
$P'(r)\varepsilon$ or $-P'(r)\varepsilon$, respectively.  At infinity its
leading Laurent term is its ordinary leading coefficient times the
corresponding power of $\varepsilon^{-1}$.  These observations justify every
``approach from a side'' and ``sign at infinity'' used below.  Passing to a
real closure of $R((\varepsilon))$ does not change any of these signs.

\begin{lemma}[Legendre congruences]\label{lem:legendre-congruences}
Let $A,B,C$ be nonzero, squarefree, and pairwise coprime.  If no ordering of
$F$ gives $A,B,C$ one common strict sign, then
\[
\begin{array}{ll}
 -BC\text{ is a square modulo }\pi &(\pi\mid A),\\
 -AC\text{ is a square modulo }\pi &(\pi\mid B),\\
 -AB\text{ is a square modulo }\pi &(\pi\mid C)
\end{array}                                             \tag{6}
\]
for every monic irreducible $\pi$.  Hence the same three assertions hold
modulo $A,B,C$, respectively.
\end{lemma}
\begin{proof}
Consider $\pi\mid A$; the other cases are cyclic.  Pairwise coprimality makes
$BC$ nonzero modulo $\pi$.

If $\pi=X-r$ is linear, then ``$-BC$ is not a square modulo $\pi$'' means
$-B(r)C(r)<0$, because over the real closed field $R$ the nonzero squares are
exactly the positive elements.  Hence $B(r)$ and $C(r)$ have the same sign.
They keep that sign sufficiently close to $r$, whereas squarefreeness makes
$A$ have a simple root and opposite signs on the two sides.  Choose the side
on which the sign of $A$ agrees with those of $B$ and $C$.  The Laurent
ordering corresponding to approaching $r$ from that side then makes all
three polynomials have one strict sign, a contradiction.

If $\pi$ is nonlinear, real closedness implies $\deg\pi=2$.  The field
$R[X]/(\pi)$ is the algebraic closure $R(\sqrt{-1})$, so every nonzero element
of it is a square.  Thus (6) holds in this case as well.

Finally, a squarefree polynomial is a product of distinct irreducibles.
The Chinese remainder theorem combines the square roots modulo those factors
and gives a square root modulo the whole polynomial.
\end{proof}

\subsection{The polynomial pivot descent}

The next lemma is the precise termination argument behind the Lean ternary
proof.  ``Real split'' means that the symbol splits after every homomorphism
to a real closed ordered field.

\begin{lemma}[Cyclic pivot step]\label{lem:pivot}
Let $A,B,C$ be nonzero squarefree pairwise-coprime polynomials satisfying the
whole-modulus congruences (6), and suppose that $\mathcal Q(A,B,C)$ is real
split.  If at least one coefficient is nonconstant, then either
$\mathcal Q(A,B,C)$ already splits over $F$, or there are nonzero squarefree
pairwise-coprime $A',B',C'$ such that
\[
\begin{split}
 \mathcal Q(A,B,C)\text{ splits}
   &\Longleftrightarrow \mathcal Q(A',B',C')\text{ splits},\\
 \mathcal Q(A,B,C)\text{ is real split}
   &\Longleftrightarrow \mathcal Q(A',B',C')\text{ is real split},\\
 \deg A'+\deg B'+\deg C'
   &<\deg A+\deg B+\deg C.
\end{split}                                             \tag{7}
\]
and the new triple again satisfies (6).
\end{lemma}
\begin{proof}
We first choose a cyclic ordering, denoted again by $(A,B,C)$, with $A$ as
pivot.  Unless all three positive degrees are equal, a coefficient of maximal
degree can be chosen so that
\[
  \deg B+\deg C<2\deg A.                                \tag{8}
\]
Choose, using (6), a representative $x$ of degree $<\deg A$ satisfying
$x^2\equiv-BC\pmod A$, and define $d$ by
\[
  x^2+BC=Ad.                                             \tag{9}
\]
The bounds on $x$ and (8) give $\deg d<\deg A$, unless $d=0$.

If $d=0$, then $x^2=-BC$, and $(0,C,x)$ is a nonzero isotropic
vector for $\langle A,B,C\rangle$.  This is the terminal branch.

Assume $d\ne0$.  Cyclic permutation identifies $\mathcal Q(A,B,C)$ with
$(-BC,-AB)$ up to splitting.  Apply (5) with $u=-BC$ and norm
$x^2+BC=Ad\ne0$.  The second slot becomes $-A^2Bd$; removing the square
factor $A^2$ gives $(-BC,-Bd)=\mathcal Q(B,C,d)$.  Thus
\[
  \mathcal Q(A,B,C)\text{ splits}
  \quad\Longleftrightarrow\quad
  \mathcal Q(B,C,d)\text{ splits}.                       \tag{10}
\]
All these changes remain valid after any field embedding, since their
nonzero denominators and norm factors remain nonzero.

Apply the parity normalization described above to the triple
$(B,C,d)$.  At each irreducible factor, remove even valuations; if the factor
occurs in two or three slots, multiply all slots by it and remove the newly
created squares.  The number of occurrences then changes from $2$ to $1$, or
from $3$ to $0$.  Hence this process produces a squarefree pairwise-coprime
triple $(A',B',C')$ without increasing total degree.  Common scaling and
individual square substitutions preserve both ordinary splitting and real
splitting, so (10) gives the first two equivalences in (7).

The new triple cannot have one common strict sign in any ordering.  Otherwise
extend that ordering to a real closure.  Both slots of its cleared symbol
would then be negative, so (3) would make the symbol nonsplit, contradicting
real splitting.  Lemma~\ref{lem:legendre-congruences}, followed by the Chinese
remainder theorem, therefore restores the whole-modulus form of (6).  The
Lean proof additionally records the direct transport through (9): a
Bézout identity for each coprime pair turns an old square witness into the
corresponding new one before the same normalization.

The degree bookkeeping is also explicit.  Before square removal the new total
degree is bounded by
\[
  \deg B+\deg C+\deg d
  <\deg A+\deg B+\deg C.
\]
Removing squares and common factors cannot increase it, proving (7).

It remains to treat
$\deg A=\deg B=\deg C=m>0$.  The orderings at $+\infty$ and $-\infty$ are
realized by Laurent series.  Real splitting therefore rules out three leading
coefficients of one sign.  After a cyclic permutation we may assume
$\operatorname{lc}(B)\operatorname{lc}(C)<0$.  Choose
$r\in R$ with
\[
  r^2=-\operatorname{lc}(B)\operatorname{lc}(C),
\]
which is possible because $R$ is real closed.  Start with the reduced square
root $x_0$ modulo $A$ and replace it by
$x=x_0+tA$, where
$t=r/\operatorname{lc}(A)$.  This does not alter the congruence modulo $A$,
but it cancels the coefficient of degree $2m$ in $x^2+BC$.  Thus (9) again
has $d=0$ or $\deg d<m$, and the preceding argument applies.
\end{proof}

\begin{proposition}[Polynomial quaternion--Legendre criterion]\label{prop:legendre}
Let $A,B,C\in R[X]$ be nonzero, squarefree, and pairwise coprime.  Suppose
that the ternary form $\langle A,B,C\rangle$ has no ordering obstruction.
Then there are polynomials $x,y,z$, not all zero, such that
\[
  Ax^2+By^2+Cz^2=0.
\]
\end{proposition}

\begin{proof}
By Lemma~\ref{lem:legendre-congruences}, the triple satisfies (6).  At every
real closed image, the three coefficients are not all of one sign, so at
least one of $-AB$ and $-AC$ is positive.  A positive element of a real
closed field is a square, and (3) shows that $\mathcal Q(A,B,C)$ is real
split.

We prove that it splits over $F$ by strong induction on
$N=\deg A+\deg B+\deg C$.  If $N=0$, the coefficients are nonzero constants.
They are not all of one sign, so one of $-AB$ and $-AC$ is positive and
hence a square in $R$; the symbol splits by (3).  If $N>0$, apply
Lemma~\ref{lem:pivot}.  The terminal alternative is finished.  Otherwise the
successor has smaller total degree, still satisfies (6), and is real split.
The induction hypothesis splits the successor, and the first equivalence in
(7) splits the original symbol.

By (3), splitting $\mathcal Q(A,B,C)=(-AB,-AC)$ gives a nonzero triple
$(x,y,z)$ with
\[
 x^2+AB\,y^2+AC\,z^2=0.
\]
Dividing by $A$ rewrites this as
\[
 A(x/A)^2+B y^2+C z^2=0.
\]
Thus $(x/A,y,z)$ is a nonzero rational-function isotropic vector.  Clearing
its common denominator gives polynomials, not all zero, satisfying
$Ax^2+By^2+Cz^2=0$.
\end{proof}

\begin{proof}[Proof of Theorem~\ref{thm:ternary}]
One direction is scalar extension.  Conversely, local isotropy and
Lemma~\ref{lem:rcdiag}, applied through the ordered real closure of
Proposition~\ref{prop:orc}, show that no ordering makes the three coefficients
have the same strict sign.  Apply the squarefree pairwise-coprime
normalization above and identify the ternary equation with the splitting
equation for the quaternion symbol
\[
  \left(-\frac ba,-\frac ca\right).
\]
Proposition~\ref{prop:legendre} gives a nonzero polynomial isotropic vector;
undoing the common scaling and the individual square substitutions gives a
nonzero isotropic vector for $\langle a,b,c\rangle$ over $F$.
\end{proof}

\section{The sign-separating compression coefficient}

We now pass from dimension $n$ to dimension $n+1$.  Split the coefficient
family into a binary head $h=(h_0,h_1)$ and a nonempty tail
$t=(t_j)_{j\in J}$.

\begin{definition}[Required sign]\label{def:reqsign}
At an ordered field point, the required sign of an auxiliary coefficient
$s$ is negative precisely when
\[
  (h_0>0\ \text{and}\ h_1>0)
  \quad\text{or}\quad
  (t_j<0\ \text{for all }j).
\]
It is positive otherwise.
\end{definition}

\begin{lemma}[Compression sign lemma]\label{lem:compress-sign}
Assume the combined family $h\sqcup t$ does not have one common strict sign.
If $s\ne0$ has the required sign, then neither
\[
  \langle h_0,h_1,s\rangle,
  \qquad
  \langle (t_j)_{j\in J},-s\rangle
\]
has one common strict sign.
\end{lemma}
\begin{proof}
If both head coefficients are positive, $s<0$.  If both are negative, the
whole family is not all negative, so the tail is not all negative and the
definition gives $s>0$.  If the head has a zero coefficient or both signs,
it already has no common strict sign.  Thus the first family has no common
strict sign.  If the tail
is all negative, then $s<0$, so $-s>0$.  If the tail is all positive, the
head cannot be all positive, hence $s>0$ and $-s<0$.  A tail already
containing a zero coefficient or both signs needs no further argument.
\end{proof}

The important point is that one rational function $s$ realizes this sign
simultaneously at every ordering.

\begin{lemma}[Squarefree polynomial representative]\label{lem:squarefree-normal}
Every $f\in F^\times$ can be written
\[
  f=Aq^2
\]
with $A\in R[X]$ nonzero and squarefree and $q\in F^\times$.
\end{lemma}
\begin{proof}
Write $f=p/r$ with nonzero $p,r\in R[X]$.  Then
$f=(pr)/r^2$.  Unique factorization writes
$pr=A h^2$, where $A$ is the product of the nonzero constant unit and
precisely those monic irreducible factors that occur to odd multiplicity.
Hence $A$ is squarefree and
$f=A(h/r)^2$.
\end{proof}

\begin{proposition}[Signed-root skeleton]\label{prop:skeleton}
Suppose $h\sqcup t$ is a regular coefficient family over $F$ with no
common strict-sign ordering.  There exists $s\in F^\times$ that has the
required sign of Definition~\ref{def:reqsign} in every real closed ordered
extension of $F$.
\end{proposition}
\begin{proof}
Use Lemma~\ref{lem:squarefree-normal} to write every coefficient in the form
\[
  a_i=A_i q_i^2,
  \qquad A_i\in R[X]\setminus\{0\}\text{ squarefree}.
\]
Multiplication by $q_i^2$ changes neither a strict sign nor isotropy.  Let
$\Sigma\subset R$ be the finite set of all real roots of all $A_i$.
On each component of $R\setminus\Sigma$, every coefficient sign, and hence
the required sign, is constant.

For $r\in\Sigma$, put $r$ into a set $E$ exactly when the required signs
immediately to the left and right of $r$ differ.  Choose
$\varepsilon\in\{1,-1\}$ to obtain the required sign on one outer interval,
and define the skeleton polynomial
\[
  S(X)=\varepsilon\prod_{r\in E}(X-r).
\]
Crossing a root in $E$ changes the sign of $S$, and crossing a root not in
$E$ does not; induction across the ordered breakpoints proves that $S$ has
the required sign at every ordinary point outside $\Sigma$.

Now let $\varphi:F\to L$ be a homomorphism to an ordered real closed field,
and put $\xi=\varphi(X)$.  Since $X-r\ne0$ in $F$, injectivity gives
$\xi\ne\varphi(r)$ for every $r\in\Sigma$.  Choose $y\in R$ having the
same finite cut with respect to $\Sigma$ as $\xi$: between two consecutive
points take their midpoint, and on an outer interval take a point one unit
beyond the endpoint.  If $\Sigma$ is empty, take $y=0$.  A field
homomorphism between real closed ordered fields
preserves order, because every positive element is a nonzero square.
Polynomial signs at $y$ and at $\xi$ therefore agree.  Indeed, over a real
closed field a polynomial is a product of real linear factors, irreducible
quadratic factors having no root and constant sign, and its leading
coefficient.  The linear-factor signs are exactly the finite-cut data.
Thus $\varphi(S)$ has the required sign.  Taking $s=S$ in $F$ proves the
claim.
\end{proof}

\section{Dimension induction}

We need one elementary reconstruction lemma.

\begin{lemma}[Binary--tail reconstruction]\label{lem:reconstruct}
Let $J$ be finite, $s\ne0$, and let $\psi=(\psi_j)_{j\in J}$.  If both
\[
  \langle a,b,s\rangle
  \quad\text{and}\quad
  \langle (\psi_j)_{j\in J},-s\rangle
\]
are isotropic, then $\langle a,b,(\psi_j)_{j\in J}\rangle$ is isotropic.
\end{lemma}
\begin{proof}
Take an isotropic vector $(y,t)$ for the second form, so
\[
  \sum_j\psi_jy_j^2=s t^2.
\]
If $t=0$, the tail itself is isotropic and we are done.  If $t\ne0$, take
an isotropic vector $(u,v,w)$ for $\langle a,b,s\rangle$.  Substituting
$s=\sum_j\psi_j(y_j/t)^2$ gives
\[
  au^2+bv^2+\sum_j\psi_j\left(\frac{w y_j}{t}\right)^2=0.
\]
The resulting vector is nonzero; otherwise $u=v=0$, and the first isotropic
vector would force $w=0$ because $s\ne0$.
\end{proof}

\begin{proposition}[Induction step]\label{prop:step}
Assume the ordering-obstruction criterion implies isotropy for diagonal forms
indexed by $J\sqcup\{*\}$, where $J\ne\varnothing$.  Then it also implies
isotropy for forms indexed by $\{0,1\}\sqcup J$.
\end{proposition}
\begin{proof}
Let $h\sqcup t$ be regular and have no common strict-sign ordering.  Choose
the coefficient $s$ from Proposition~\ref{prop:skeleton}.  In every real
closed ordered extension, restriction of the order to $F$ and the original
obstruction hypothesis let Lemma~\ref{lem:compress-sign} apply.
Lemma~\ref{lem:rcdiag} then makes both compressed forms isotropic.

The ternary local--global theorem therefore makes
$\langle h_0,h_1,s\rangle$ isotropic over $F$.  The augmented tail has no
ordering obstruction: extend any proposed definite ordering to a real closure
by Proposition~\ref{prop:orc}.  There $s$ has the required sign, so
Lemma~\ref{lem:compress-sign} contradicts definiteness of $\langle t,-s\rangle$.
The induction hypothesis therefore makes
$\langle t,-s\rangle$ isotropic over $F$.  Apply
Lemma~\ref{lem:reconstruct}.
\end{proof}

\begin{proof}[Proof of Theorem~\ref{thm:diagonal}]
Lemma~\ref{lem:definite} gives (i)$\Rightarrow$(ii).  The ternary case of
(ii)$\Rightarrow$(i) is Theorem~\ref{thm:ternary}, because every real closed
image is indefinite and hence isotropic by Lemma~\ref{lem:rcdiag}.
Proposition~\ref{prop:step} now proves the implication in every dimension
$n\ge3$ by induction.  Reindexing by a bijection with $\{0,\ldots,n-1\}$
handles every finite index type.  If a coefficient is zero, its coordinate
vector handles the omitted nonregular case.

The implication (i)$\Rightarrow$(iii) is scalar extension.  Conversely,
apply (iii) to the ordered real closure extending any proposed definite
ordering.  Lemma~\ref{lem:rcdiag} gives a contradiction, so (ii) holds and the
already proved implication gives (i).
\end{proof}

\section{General quadratic forms}

\begin{lemma}[Orthogonal diagonalization]\label{lem:diagonalization}
Over a field of characteristic different from $2$, every quadratic form on a
finite-dimensional vector space has an orthogonal basis.
\end{lemma}
\begin{proof}
Let $B_q(x,y)=q(x+y)-q(x)-q(y)$ be the associated symmetric bilinear form, so
$B_q(v,v)=2q(v)$.  We argue by induction on the dimension.  If $q=0$, every
basis is orthogonal.  Otherwise choose $v$ with $q(v)\ne0$.  For every $x$ put
\[
  x^\perp=x-\frac{B_q(x,v)}{2q(v)}v.
\]
Then $B_q(x^\perp,v)=0$, and every $x$ is uniquely a sum of a multiple of
$v$ and an element of $v^\perp$.  Hence
$V=Fv\oplus v^\perp$.  Apply the induction hypothesis to the restriction of
$q$ to $v^\perp$ and adjoin $v$ to the resulting orthogonal basis.
\end{proof}

\begin{proof}[Proof of Theorem~\ref{thm:qf}]
The characteristic is zero, so Lemma~\ref{lem:diagonalization} gives an
orthogonal basis $e_1,\ldots,e_n$.  In that basis $q$ is isometric to
\[
  \sum_{i=1}^n q(e_i)X_i^2.
\]
This statement includes the degenerate case: some $q(e_i)$ may be zero, and
Theorem~\ref{thm:diagonal} already handles it.  The isometry remains an
isometry after scalar extension, and an invertible linear change of variables
preserves the existence of a nonzero isotropic vector.  Apply
Theorem~\ref{thm:diagonal} to the coefficient family
$(q(e_i))$ and transport both directions through these isometries.
\end{proof}

\section{Correspondence with Lean}

The principal correspondence is shown below.  A mathematical lemma may map to
several smaller Lean declarations; the table names the public consumer or the
module-level endpoint.

\begingroup
\scriptsize
\begin{center}
\begin{tabular}{p{0.31\linewidth}p{0.62\linewidth}}
Mathematical result & Lean declaration \\ \hline
Definiteness obstruction & \lean{RatFuncWittLocalGlobal.RatFunc.not_finiteSameStrictSignOrdering_of_isotropic} \\
Real-closed diagonal criterion & \lean{RatFuncWittLocalGlobal.Diagonal.not_isotropic_iff_all_pos_or_all_neg} \\
Preordering extension & \lean{RingPreordering.exists_le_isOrdering} \\
Positive projection criterion & \lean{RingPreordering.exists_orderedAlgebraExtension_of_projection} \\
Odd power-basis lemma & \lean{RatFuncWittLocalGlobal.OrderedRealClosure.minus_one_notMem_baseSquareSpan_of_odd_powerBasis} \\
Ordered real closure & \lean{RatFuncWittLocalGlobal.OrderedRealClosure.orderedFieldRealClosedExtension_same_universe} \\
Quadratic maximality step & \lean{RatFuncWittLocalGlobal.OrderedRealClosure.nonnegative_isSquare_of_isMax} \\
Odd-degree maximality step & \lean{RatFuncWittLocalGlobal.OrderedRealClosure.exists_root_of_odd_natDegree_of_isMax} \\
Ternary theorem & \lean{RatFuncWittLocalGlobal.RatFunc.ternary_isotropic_iff_forall_realClosedExtension_same_universe} \\
Finite Legendre congruences & \lean{RatFuncWittLocalGlobal.RatFunc.PolynomialLegendreLocalConditions} \\
Whole-modulus constructor & \lean{RatFuncWittLocalGlobal.polynomialTernaryFiniteTameResiduesTrivial_globalSquareMod} \\
Nonlinear residue square & \lean{RatFuncWittLocalGlobal.RatFunc.nonlinearIrreducibleDivisorSquareMod_of_realClosed} \\
Cleared-binary normalization & \lean{RatFuncWittLocalGlobal.RatFunc.normalizedTernaryLocalGlobal_iff_squarefree_nonterminal_right} \\
Pivot successor & \lean{RatFuncWittLocalGlobal.polynomialTernaryClearedQuaternionSplit_descent_wholeMod_normalized_right_pivot_complete} \\
Total-degree induction & \lean{RatFuncWittLocalGlobal.polynomialTernaryClearedQuaternionSplit_of_wholeMod_realSplit} \\
Legendre endpoint bridge & \lean{RatFuncWittLocalGlobal.polynomialLegendreConstruction_of_normalizedOrderingRealization} \\
Squarefree normal form & \lean{RatFuncWittLocalGlobal.finitePolynomialState_exists} \\
Compression sign lemma & \lean{RatFuncWittLocalGlobal.binaryTailCompressionRequiredSign_spec} \\
Skeleton transport & \lean{RatFuncWittLocalGlobal.FinitePolynomialState.binaryTailSkeleton_allOrdered} \\
Finite-cut sign transport & \lean{RatFuncWittLocalGlobal.polynomial_eval₂_pos_iff_eval_of_same_finite_cut} \\
Reconstruction & \lean{RatFuncWittLocalGlobal.Diagonal.isotropic_sumType_of_ternary_of_tailAdjoinNeg} \\
Induction step & \lean{RatFuncWittLocalGlobal.diagonalOrderingObstructionSufficiency_sum_fin_two} \\
Total diagonal theorem & \lean{RatFuncWittLocalGlobal.diagonal_isotropic_iff_forall_realClosedExtension_of_card_ge_three_same_universe_total} \\
Isometry invariance & \lean{RatFuncWittLocalGlobal.quadraticMap_isometryEquiv_not_anisotropic_iff} \\
Base-change diagonal isometry & \lean{RatFuncWittLocalGlobal.quadraticFormBaseChangeIsometryWeightedSumSquares} \\
General quadratic form & \lean{RatFuncWittLocalGlobal.quadraticForm_not_anisotropic_iff_forall_realClosedExtension_same_universe}
\end{tabular}
\end{center}
\endgroup

The reproducible audit in \texttt{scripts/AxiomAudit.lean} checks that the main
theorems use only \texttt{propext}, \texttt{Classical.choice}, and
\texttt{Quot.sound}.  No project-specific axiom is introduced.
The bundle and transport declarations in the proof implement the congruences,
normalizations, finite-cut argument, and changes of coordinates proved above;
they do not supply an additional unstated mathematical theorem.

\begin{thebibliography}{9}
\bibitem{Lam}
T. Y. Lam,
\emph{Introduction to Quadratic Forms over Fields},
Graduate Studies in Mathematics 67, American Mathematical Society, 2005.

\bibitem{Pfister}
A. Pfister,
\emph{Quadratic Forms with Applications to Algebraic Geometry and Topology},
London Mathematical Society Lecture Note Series 217,
Cambridge University Press, 1995.
\end{thebibliography}

\end{document}
